% Gauge Fields from the Gauss Equation of the Metric Bundle
% For submission to arXiv [hep-th]

\documentclass[12pt,a4paper]{article}

% ---- Packages ----
\usepackage[utf8]{inputenc}
\usepackage[T1]{fontenc}
\usepackage{amsmath,amssymb,amsthm}
\usepackage{mathtools}
\usepackage{geometry}
\usepackage{hyperref}
\usepackage{cleveref}
\usepackage{enumitem}
\usepackage{booktabs}
\usepackage{array}
\usepackage{cite}

\geometry{margin=1in}

% ---- Theorem environments ----
\newtheorem{theorem}{Theorem}[section]
\newtheorem{proposition}[theorem]{Proposition}
\newtheorem{lemma}[theorem]{Lemma}
\newtheorem{corollary}[theorem]{Corollary}
\theoremstyle{definition}
\newtheorem{definition}[theorem]{Definition}
\newtheorem{remark}[theorem]{Remark}

% ---- Custom commands ----
\newcommand{\Tr}{\operatorname{Tr}}
\newcommand{\tr}{\operatorname{tr}}
\newcommand{\Met}{\operatorname{Met}}
\newcommand{\GL}{\operatorname{GL}}
\newcommand{\SO}{\operatorname{SO}}
\newcommand{\SU}{\operatorname{SU}}
\newcommand{\U}{\operatorname{U}}
\newcommand{\Spin}{\operatorname{Spin}}
\newcommand{\Cl}{\operatorname{Cl}}
\newcommand{\ad}{\operatorname{ad}}
\newcommand{\diag}{\operatorname{diag}}
\newcommand{\Sym}{\operatorname{Sym}}
\newcommand{\R}{\mathbb{R}}
\newcommand{\C}{\mathbb{C}}
\newcommand{\Z}{\mathbb{Z}}

% ---- Title ----
\title{\textbf{Gauge Fields from the Gauss Equation\\ of the Metric Bundle}}

\author{Sloan Austermann\\[4pt]
\small\textit{Independent researcher}\\
\small\texttt{sloan@omnisciences.org}}

\date{\today}

% ============================================================
\begin{document}
\maketitle

% ---- Abstract ----
\begin{abstract}
In earlier work, we showed that the DeWitt supermetric on the space of Lorentzian metrics over a four-dimensional spacetime $X$ has signature~$(6,4)$, yielding the Pati-Salam gauge group from the maximal compact subgroup of~$\SO(6,4)$.
In this paper, we carry out the Gauss-Codazzi-Ricci decomposition of the scalar curvature of the fourteen-dimensional metric bundle $Y^{14} = \Met(X)$ restricted to a metric section $g\colon X \hookrightarrow Y$.
We show that the Einstein-Hilbert, Yang-Mills, and extrinsic curvature terms emerge with the correct relative signs: (i)~$+R_X$ for gravity, (ii)~$-|II|^2$ for the extrinsic curvature (torsion), and (iii)~$-(h/4)|F|^2$ with $h > 0$ for the Yang-Mills gauge fields.
The fibre $\GL^+(4,\R)/\SO(4)$ is a symmetric space of non-compact type with scalar curvature $R_{\mathrm{fibre}} = -36$.
We compute the gauge kinetic metric $h_{ab}$ for the $\SO(4) = \SU(2)_L \times \SU(2)_R$ gauge sector via the Kaluza-Klein mechanism and find $h_L = h_R = 6\cdot I_3$, confirming left-right symmetry $g_L = g_R$ at the unification scale and the absence of ghosts in the gauge sector.
The non-vanishing commutators of the shape operators, $[A_m, A_n] \neq 0$ for $24$ of the $45$ normal-direction pairs, provide the non-abelian field strength via the Ricci equation.
\end{abstract}

\vspace{10pt}
\noindent\textit{Keywords:} Gauss equation, metric bundle, Kaluza-Klein reduction, Yang-Mills sign, symmetric space

\newpage
\tableofcontents
\newpage

% ============================================================
\section{Introduction}
\label{sec:intro}
% ============================================================

The metric bundle programme~\cite{Paper1,Paper2} derives the Pati-Salam gauge group $\SU(4) \times \SU(2)_L \times \SU(2)_R$ from the geometry of $Y^{14} = \Met(X^4)$, the bundle of Lorentzian metrics over four-dimensional spacetime.
The DeWitt supermetric~\cite{DeWitt1967} on the ten-dimensional fibre of symmetric two-tensors has signature~$(6,4)$ when evaluated on a Lorentzian background, giving rise to a normal bundle with structure group~$\SO(6,4)$.

Having established the gauge group~\cite{Paper1} and the fermion content~\cite{Paper2}, a crucial test of the framework is whether the \emph{dynamics}---specifically the effective four-dimensional action---emerges with the correct signs.
The standard Kaluza-Klein mechanism~\cite{KaluzaKlein,Bailin1987} extracts gauge fields from the isometries of the internal space, but the metric bundle presents a non-standard setting: the fibre $\GL^+(4,\R)/\SO(4)$ is \emph{non-compact}, the DeWitt metric has \emph{indefinite} signature, and the embedding is via a \emph{section} rather than a product decomposition.

In this paper, we use the Gauss-Codazzi-Ricci equations for the submanifold $g(X) \hookrightarrow Y$ to decompose the fourteen-dimensional scalar curvature and extract the four-dimensional effective action.
The main results are:

\begin{enumerate}[nosep]
\item The second fundamental form of the trivial section has $|II|^2 = 2$ and mean curvature $|H|^2 = -1$ (the negative value arising from the indefinite DeWitt metric).
\item The Gauss equation gives $R_Y = R_X + |H|^2 - |II|^2 + 2\,\mathrm{Ric}_{\mathrm{mixed}} + R^\perp$, with $+R_X$ (correct Einstein-Hilbert sign) and $-|II|^2$ (correct sign for extrinsic curvature minimisation).
\item The fibre curvature, computed from the symmetric space formula $R(X,Y)Z = -[[X,Y],Z]$, gives scalar curvature $R_{\mathrm{fibre}} = -36$.
\item The gauge kinetic metric $h_{ab}$ for the $\SO(4) = \SU(2)_L \times \SU(2)_R$ gauge sector is positive definite on the traceless sector, giving the correct sign $-(h/4)|F|^2$ in the Yang-Mills action.
\item Left-right symmetry $h_L = h_R = 6\cdot I_3$ is exact, confirming $g_L = g_R$ at the compactification scale.
\end{enumerate}

\paragraph{Outline.}
Section~\ref{sec:fibre} reviews the fibre geometry of the symmetric space $\GL^+(4,\R)/\SO(4)$.
Section~\ref{sec:gauss} presents the Gauss equation and computes the second fundamental form.
Section~\ref{sec:fibre_curvature} computes the fibre curvature tensor from the symmetric space structure.
Section~\ref{sec:yang_mills} derives the Yang-Mills sign from the gauge kinetic metric.
Section~\ref{sec:su2_decomp} performs the decomposition under $\SU(2)_L \times \SU(2)_R$.
Section~\ref{sec:action} assembles the full effective four-dimensional action.
Section~\ref{sec:discussion} discusses implications and open questions.


% ============================================================
\section{The Fibre Geometry}
\label{sec:fibre}
% ============================================================

The fibre of the metric bundle $Y^{14}$ at a point $x \in X$ is the space of positive-definite symmetric bilinear forms on $T_xX$, which is diffeomorphic to the symmetric space $\GL^+(4,\R)/\SO(4)$.

\subsection{The Cartan decomposition}

The Lie algebra $\mathfrak{gl}(4,\R)$ admits the Cartan decomposition
\begin{equation}
\mathfrak{gl}(4,\R) = \mathfrak{k} \oplus \mathfrak{p}\,,
\end{equation}
where $\mathfrak{k} = \mathfrak{o}(4)$ (antisymmetric matrices) and $\mathfrak{p} = S^2(\R^4)$ (symmetric matrices).
The Lie bracket relations are:
\begin{equation}
[\mathfrak{k}, \mathfrak{k}] \subset \mathfrak{k}\,,\qquad
[\mathfrak{k}, \mathfrak{p}] \subset \mathfrak{p}\,,\qquad
[\mathfrak{p}, \mathfrak{p}] \subset \mathfrak{k}\,.
\end{equation}
The last relation is key: for $h, k \in \mathfrak{p} = S^2(\R^4)$, the commutator $[h,k] = hk - kh$ is antisymmetric and hence lies in $\mathfrak{k} = \mathfrak{o}(4)$.

\subsection{The DeWitt metric on $\mathfrak{p}$}

The DeWitt supermetric~\cite{DeWitt1967} provides a natural inner product on $\mathfrak{p}$:
\begin{equation}
\label{eq:dewitt}
\mathcal{G}(h,k) = \tr(hk) - \tfrac{1}{2}\tr(h)\tr(k)\,.
\end{equation}
In the orthonormal basis $\{e_m\}_{m=1}^{10}$ for $S^2(\R^4)$ under the trace inner product $\langle h, k \rangle = \tr(hk)$---with diagonal elements $e_{(i,i)} = E_{ii}$ and off-diagonal elements $e_{(i,j)} = (E_{ij} + E_{ji})/\sqrt{2}$ for $i < j$---the DeWitt metric matrix~$G_{mn}$ has eigenvalues:
\begin{equation}
\underbrace{+1, \ldots, +1}_{9}, \quad -1\,.
\end{equation}
The negative direction is the conformal (trace) mode $\propto \sum_i E_{ii}$.
The traceless subspace $S^2_0(\R^4)$ has positive-definite DeWitt metric.

\begin{remark}
The signature~$(9,1)$ applies to the Euclidean background $g = \delta$.
For the Lorentzian background $g = \diag(-1,1,1,1)$, the signature changes to~$(6,4)$ as shown in~\cite{Paper1}.
In this paper, we work with the Euclidean fibre to compute the intrinsic fibre curvature and gauge structure, which is the relevant computation for the structure group of the normal bundle at the identity coset.
\end{remark}


% ============================================================
\section{The Gauss Equation}
\label{sec:gauss}
% ============================================================

\subsection{Setup}

A metric section $g\colon X^4 \to Y^{14}$ embeds $X$ as a four-dimensional submanifold of the fourteen-dimensional total space.
The Gauss equation~\cite{KobayashiNomizu} relates the curvatures:
\begin{equation}
\label{eq:gauss}
\langle R_Y(u,v)w, z \rangle = \langle R_X(u,v)w, z \rangle + \langle II(u,w), II(v,z) \rangle - \langle II(u,z), II(v,w) \rangle\,,
\end{equation}
where $u, v, w, z \in T_xX$ and $II$ is the second fundamental form of the embedding.

Taking traces, the scalar curvature decomposes as:
\begin{equation}
\label{eq:scalar_decomp}
R_Y\big|_{g(X)} = R_X + |H|^2 - |II|^2 + 2\,\mathrm{Ric}_{\mathrm{mixed}} + R^\perp\,,
\end{equation}
where $H = \tr_g II$ is the mean curvature vector, $\mathrm{Ric}_{\mathrm{mixed}}$ collects the tangent-normal sectional curvatures, and $R^\perp$ is the contribution from the normal-normal sectional curvatures.

\subsection{The second fundamental form}

For the trivial section $g = \bar{g}$ (constant background metric) over a flat base, the second fundamental form reduces to the mixed Christoffel symbols of the chimeric metric:
\begin{equation}
II^m_{\mu\nu}\big|_{g=\bar{g}} = \Gamma^m_{\mu\nu}\big|_Y = \tfrac{1}{2}\, G^{mk}\, \frac{\partial g_{\mu\nu}}{\partial y^k}\bigg|_{y=0} = \tfrac{1}{2}\, G^{mk}\, e^k_{(\mu\nu)}\,,
\end{equation}
where $e^k_{(\mu\nu)}$ is the $(\mu,\nu)$ component of the $k$-th basis element of $S^2(\R^4)$, and $G^{mk}$ is the inverse DeWitt metric.

\begin{proposition}[Second fundamental form invariants]
\label{prop:II}
At the trivial section over a flat base:
\begin{align}
|II|^2 &= G_{mn}\, g^{\mu\rho} g^{\nu\sigma}\, II^m_{\mu\nu}\, II^n_{\rho\sigma} = 2\,, \\
|H|^2 &= G_{mn}\, H^m\, H^n = -1\,,
\end{align}
where $H^m = g^{\mu\nu}\, II^m_{\mu\nu}$ is the mean curvature vector.
\end{proposition}

\begin{proof}
Direct computation using the inverse DeWitt metric $G^{mk}$ and the basis elements $e^k_{(\mu\nu)}$.
The mean curvature vector $H^m = \sum_\mu \Gamma^m_{\mu\mu}$ has nonzero components only along the diagonal basis elements.
Its norm-squared $|H|^2 = G_{mn} H^m H^n = -1$ is negative because $H$ points predominantly along the conformal (trace) direction, which has negative DeWitt norm.
\end{proof}

\begin{remark}
The decomposition $II = \frac{1}{d}H \otimes g + II_0$ separates the trace ($H$) from the traceless ($II_0$) part, giving $|II|^2 = \frac{1}{d}|H|^2 + |II_0|^2$.
With $|II|^2 = 2$ and $|H|^2/d = -1/4$, the traceless contribution is $|II_0|^2 = 2 + 1/4 = 9/4$.
\end{remark}


% ============================================================
\section{Fibre Curvature}
\label{sec:fibre_curvature}
% ============================================================

\subsection{The symmetric space curvature formula}

For a symmetric space $G/K$ with Cartan decomposition $\mathfrak{g} = \mathfrak{k} \oplus \mathfrak{p}$, the Riemannian curvature tensor at the identity coset is determined algebraically~\cite{Helgason1978}:
\begin{equation}
\label{eq:symspace_curvature}
R(X,Y)Z = -[[X,Y]_\mathfrak{k}, Z]\,,\qquad X, Y, Z \in \mathfrak{p}\,.
\end{equation}
For $\GL^+(4,\R)/\SO(4)$, with $X, Y \in \mathfrak{p} = S^2(\R^4)$ being symmetric matrices, we have $[X,Y] = XY - YX \in \mathfrak{k} = \mathfrak{o}(4)$ (already antisymmetric), so $[X,Y]_\mathfrak{k} = [X,Y]$ and
\begin{equation}
R(X,Y)Z = -[XY - YX, Z] = -(XY - YX)Z + Z(XY - YX)\,.
\end{equation}

\subsection{Riemann tensor and curvature invariants}

\begin{proposition}[Fibre curvature invariants]
\label{prop:fibre_curv}
The symmetric space $\GL^+(4,\R)/\SO(4)$ with the DeWitt metric has:
\begin{enumerate}[nosep]
\item Riemann tensor $R^m{}_{nrs}$ satisfying all standard symmetries: $R_{mnrs} = -R_{mnsr} = -R_{nmrs} = R_{rsmn}$.
\item All sectional curvatures $K(e_m, e_n) \leq 0$ (non-compact type).
\item Scalar curvature $R_{\mathrm{fibre}} = -36$.
\end{enumerate}
\end{proposition}

\begin{proof}
The Riemann tensor components are computed from~\eqref{eq:symspace_curvature}:
$R_{mnrs} = \mathcal{G}(R(e_r, e_s)e_n, e_m)$.
The symmetries follow from the general theory of symmetric spaces.
The non-positivity of sectional curvatures is a standard result for symmetric spaces of non-compact type~\cite{Helgason1978}: the Killing form is negative definite on $\mathfrak{k}$, implying $K(X,Y) = -|[X,Y]|^2/(|X|^2|Y|^2 - \langle X,Y\rangle^2) \leq 0$.
The scalar curvature $R_{\mathrm{fibre}} = G^{mn} G^{rs} R_{mrns} = -36$ is obtained by explicit contraction.
\end{proof}


% ============================================================
\section{The Yang-Mills Sign}
\label{sec:yang_mills}
% ============================================================

The most critical test of any Kaluza-Klein framework is whether the Yang-Mills kinetic term emerges with the correct sign~\cite{Bailin1987}.
A wrong sign would indicate ghost gauge fields, rendering the theory non-unitary.

\subsection{Shape operators and non-abelian structure}

The shape operator (Weingarten map) for a normal vector $\xi^m e_m$ is:
\begin{equation}
A_\xi\colon T_xX \to T_xX\,,\qquad \langle A_\xi(u), v \rangle = \langle II(u,v), \xi \rangle\,.
\end{equation}
At the trivial section, $A_m = \Gamma_m$ (the mixed Christoffel symbol matrix viewed as a $4 \times 4$ endomorphism of $T_xX$).

\begin{proposition}[Non-abelian shape operators]
\label{prop:shape}
Of the $\binom{10}{2} = 45$ pairs of normal directions, exactly $24$ have non-vanishing commutators $[A_m, A_n] \neq 0$, with
\begin{equation}
\sum_{m < n} |[A_m, A_n]|^2 = \tfrac{9}{8}\,.
\end{equation}
\end{proposition}

The non-commutativity of shape operators is the geometric origin of the non-abelian gauge field strength.
By the Ricci equation,
\begin{equation}
\label{eq:ricci}
\langle R^\perp(u,v)\xi, \eta \rangle = \langle R_Y(u,v)\xi, \eta \rangle + \langle [A_\xi, A_\eta]u, v \rangle\,,
\end{equation}
the normal curvature $R^\perp$ receives a contribution from the commutator $[A_\xi, A_\eta]$, which for a gauge field perturbation gives the non-abelian part of the field strength $F = dA + A \wedge A$.

\subsection{The gauge kinetic metric}

The gauge algebra at the identity coset is $\mathfrak{k} = \mathfrak{o}(4)$, acting on the fibre $\mathfrak{p} = S^2(\R^4)$ via the adjoint representation:
\begin{equation}
\ad_k(h) = [k, h] = kh - hk\,,\qquad k \in \mathfrak{o}(4),\; h \in S^2(\R^4)\,.
\end{equation}

In the Kaluza-Klein reduction~\cite{Bailin1987}, the Yang-Mills kinetic term takes the form
\begin{equation}
\label{eq:YM_action}
S_{\mathrm{YM}} = -\tfrac{1}{4}\, h_{ab}\! \int_X F^a_{\mu\nu} F^{b\,\mu\nu}\, \sqrt{|g|}\, d^4x\,,
\end{equation}
where the gauge kinetic metric $h_{ab}$ is determined by the DeWitt inner product on the gauge orbits:
\begin{equation}
\label{eq:gauge_kinetic}
h_{ab} = \sum_{m,n} G^{\mathrm{DW}}_{mn}\, (\ad_a)^m{}_r\, (\ad_b)^n{}_r = \tr\!\big((\ad_a)^T\! \cdot G^{\mathrm{DW}} \cdot \ad_b\big)\,,
\end{equation}
where $(\ad_a)^m{}_n$ denotes the matrix of $\ad_{k_a}$ in the basis $\{e_m\}$ of $\mathfrak{p}$.

\begin{theorem}[Correct Yang-Mills sign]
\label{thm:YM_sign}
The gauge kinetic metric $h_{ab}$ is positive definite on the gauge algebra $\mathfrak{o}(4)$:
\begin{equation}
h_{ab} > 0\qquad \text{for all nonzero } k_a \in \mathfrak{o}(4)\,.
\end{equation}
Consequently, the Yang-Mills action~\eqref{eq:YM_action} has the correct negative sign, and the gauge sector contains no ghosts.
\end{theorem}

\begin{proof}
The key observation is that the gauge transformations $\ad_k\colon S^2(\R^4) \to S^2(\R^4)$ preserve the trace: $\tr([k,h]) = \tr(kh - hk) = 0$ for all $h$.
Therefore the image of $\ad_k$ lies entirely in the \emph{traceless} subspace $S^2_0(\R^4)$, on which the DeWitt metric is positive definite.

Since $G^{\mathrm{DW}}_{mn}$ restricted to traceless symmetric matrices is a positive-definite form, and the gauge field components live entirely in this subspace, the contraction $h_{ab} = \tr((\ad_a)^T\! \cdot G^{\mathrm{DW}} \cdot \ad_b)$ is positive semi-definite.
It is positive definite on $\mathfrak{o}(4)$ because each $\ad_{k_a}$ has nonzero image in $S^2_0(\R^4)$ (no element of $\mathfrak{o}(4)$ centralises all of $S^2(\R^4)$).
\end{proof}

\begin{remark}
This is a non-trivial result.
The DeWitt metric has signature~$(9,1)$ on the full fibre, and one might worry that the indefinite signature could produce a wrong-sign gauge kinetic term.
The saving feature is that gauge transformations are volume-preserving (traceless), and the trace direction is precisely the one negative direction.
\end{remark}


% ============================================================
\section{Decomposition under $\SU(2)_L \times \SU(2)_R$}
\label{sec:su2_decomp}
% ============================================================

The gauge group $\SO(4) \cong [\SU(2)_L \times \SU(2)_R]/\Z_2$ decomposes via the self-dual/anti-self-dual splitting of $\mathfrak{o}(4)$.

\subsection{Self-dual and anti-self-dual generators}

In four dimensions, the Hodge star on 2-forms gives the decomposition $\Lambda^2(\R^4) = \Lambda^2_+ \oplus \Lambda^2_-$.
Correspondingly, $\mathfrak{o}(4) = \mathfrak{su}(2)_L \oplus \mathfrak{su}(2)_R$ with generators:
\begin{align}
\label{eq:su2_generators}
\mathfrak{su}(2)_L&\colon\quad L_1 = \tfrac{1}{2}(e_{12} + e_{34}),\; L_2 = \tfrac{1}{2}(e_{13} - e_{24}),\; L_3 = \tfrac{1}{2}(e_{14} + e_{23})\,,\\
\mathfrak{su}(2)_R&\colon\quad R_1 = \tfrac{1}{2}(e_{12} - e_{34}),\; R_2 = \tfrac{1}{2}(e_{13} + e_{24}),\; R_3 = \tfrac{1}{2}(e_{14} - e_{23})\,,
\end{align}
where $e_{ij}$ is the antisymmetric matrix with $(e_{ij})_{ij} = +1$, $(e_{ij})_{ji} = -1$.
These satisfy $[L_i, L_j] = \epsilon_{ijk} L_k$, $[R_i, R_j] = \epsilon_{ijk} R_k$, and $[L_i, R_j] = 0$.

\subsection{The representation on $S^2(\R^4)$}

Under $\SU(2)_L \times \SU(2)_R$, the fundamental representation $\R^4 \cong (\mathbf{2}, \mathbf{2})$ and hence:
\begin{equation}
S^2(\R^4) = S^2(\mathbf{2},\mathbf{2}) = (\mathbf{3},\mathbf{3}) \oplus (\mathbf{1},\mathbf{1})\,.
\end{equation}
The singlet $(\mathbf{1},\mathbf{1})$ is the trace (conformal mode), and the nine-dimensional traceless subspace carries the $(\mathbf{3},\mathbf{3})$ representation.

\begin{proposition}[Equal gauge couplings]
\label{prop:equal_couplings}
The gauge kinetic metric $h_{ab}$ is block-diagonal under the $\SU(2)_L \oplus \SU(2)_R$ decomposition:
\begin{equation}
h_L = 6\cdot I_3\,,\qquad h_R = 6\cdot I_3\,,
\end{equation}
and consequently $g_L/g_R = \sqrt{h_R/h_L} = 1$.
Left-right symmetry holds exactly at the compactification scale.
\end{proposition}

\begin{proof}
The $\SU(2)_L$ generators $\{L_i\}$ and $\SU(2)_R$ generators $\{R_i\}$ commute, so $h_{L_i R_j} = 0$ by the orthogonality of the self-dual and anti-self-dual subspaces under the DeWitt metric.
The diagonal entries are computed from~\eqref{eq:gauge_kinetic}:
\begin{equation}
(h_L)_{ij} = \tr\!\big((\ad_{L_i})^T\! \cdot G^{\mathrm{DW}} \cdot \ad_{L_j}\big) = 6\,\delta_{ij}\,,
\end{equation}
and similarly for $h_R$.
The equality $h_L = h_R$ follows from the parity symmetry of the DeWitt metric: the metric~\eqref{eq:dewitt} is invariant under the exchange of any two spatial indices, which maps $L_i \leftrightarrow R_i$ (up to signs absorbed by the $I_3$ factor).
\end{proof}

\subsection{Casimir operators}

The quadratic Casimirs $C_2(\SU(2)_L)$ and $C_2(\SU(2)_R)$ on $S^2(\R^4)$ take eigenvalues:
\begin{equation}
C_2^{L,R}\big|_{(\mathbf{3},\mathbf{3})} = j(j+1)\big|_{j=1} = 2\,,\qquad
C_2^{L,R}\big|_{(\mathbf{1},\mathbf{1})} = 0\,,
\end{equation}
confirming the representation-theoretic decomposition.


% ============================================================
\section{The Effective Four-Dimensional Action}
\label{sec:action}
% ============================================================

Assembling all contributions from the Gauss-Codazzi-Ricci decomposition, the scalar curvature of $Y$ restricted to the section $g(X)$ decomposes as:

\begin{equation}
\label{eq:full_decomp}
R_Y\big|_{g(X)} = R_X + |H|^2 - |II|^2 + 2\,\mathrm{Ric}_{\mathrm{mixed}} + R^\perp\,.
\end{equation}

\subsection{Sign structure}

We summarise the physical interpretation of each term:

\begin{center}
\renewcommand{\arraystretch}{1.3}
\begin{tabular}{lll}
\toprule
\textbf{Geometric term} & \textbf{Sign in $R_Y$} & \textbf{Physical identification} \\
\midrule
$R_X$ & $+$ & Einstein-Hilbert (gravity) \\
$|H|^2$ & $+$ & Conformal mode (non-propagating) \\
$-|II|^2$ & $-$ & Extrinsic curvature / torsion \\
$R^\perp$ & & Yang-Mills gauge fields \\
$\mathrm{Ric}_{\mathrm{mixed}}$ & & Gravity-gauge coupling \\
\bottomrule
\end{tabular}
\end{center}

The action on $Y$ is $S_Y = \frac{1}{16\pi G_{14}} \int_Y R_Y\, \mathrm{vol}_Y$.
Restricted to the section:
\begin{equation}
S_{\mathrm{eff}} = \frac{V_{\mathrm{eff}}}{16\pi G_{14}} \int_X \Big[ R_X + |H|^2 - |II|^2 + 2\,\mathrm{Ric}_{\mathrm{mixed}} + R^\perp \Big] \mathrm{vol}_X\,,
\end{equation}
where $V_{\mathrm{eff}}$ is an effective localisation factor from the normal directions.

\subsection{Identification with Standard Model terms}

\begin{enumerate}
\item \textbf{Gravity:} $+R_X$ with positive coefficient gives the Einstein-Hilbert action with the correct sign.

\item \textbf{Torsion/extrinsic curvature:} $-|II|^2$ enters with a minus sign.
For a general section $g(x) = \bar{g} + h(x)$, the second fundamental form encodes the deviation of the section from being totally geodesic.
The identification $|II|^2 \sim |T|^2$ (torsion squared) gives a bounded-below torsion action.

\item \textbf{Yang-Mills:} The gauge kinetic contribution comes from the dynamical part of the normal curvature.
For gauge field perturbations $A^a_\mu$ along the Killing directions of $\mathfrak{o}(4)$, the Yang-Mills term is
\begin{equation}
S_{\mathrm{YM}} = -\frac{V_{\mathrm{eff}}}{64\pi G_{14}}\, h_{ab} \int_X F^a_{\mu\nu} F^{b\,\mu\nu}\, \mathrm{vol}_X\,,
\end{equation}
with $h_{ab}$ positive definite (Theorem~\ref{thm:YM_sign}), giving the standard Yang-Mills action with correct sign.

\item \textbf{Conformal mode:} $+|H|^2 = -1$ at the trivial section.
The positive coefficient in the action with negative value of $|H|^2$ reflects the well-known conformal factor problem of quantum gravity~\cite{Giulini2009}.
In the metric bundle framework, the mean curvature $H$ is related to the volume-stabilisation potential and is constrained by the section condition, not a freely propagating field.
\end{enumerate}

\subsection{The three sign tests}

The viability of the metric bundle framework depends on three sign tests, all of which are passed:

\begin{center}
\renewcommand{\arraystretch}{1.3}
\begin{tabular}{llc}
\toprule
\textbf{Test} & \textbf{Requirement} & \textbf{Result} \\
\midrule
Einstein-Hilbert & $+R_X$ in $S_{\mathrm{eff}}$ & \checkmark \\
Torsion/extrinsic curvature & $-|II|^2$ (bounded below) & \checkmark \\
Yang-Mills & $-\frac{h}{4}|F|^2$ with $h > 0$ & \checkmark \\
\bottomrule
\end{tabular}
\end{center}


% ============================================================
\section{Discussion}
\label{sec:discussion}
% ============================================================

We have shown that the Gauss-Codazzi-Ricci decomposition of the metric bundle $Y^{14} = \Met(X^4)$ produces an effective four-dimensional action with the correct sign structure for gravity, gauge fields, and extrinsic curvature.
The key results are:

\begin{enumerate}
\item The \emph{non-compact} fibre $\GL^+(4,\R)/\SO(4)$ does not produce ghosts because gauge transformations are volume-preserving, and the DeWitt metric is positive definite on the traceless (volume-preserving) subspace.

\item Left-right symmetry $g_L = g_R$ is an \emph{exact} consequence of the parity invariance of the DeWitt metric, rather than an assumption.

\item The fibre scalar curvature $R_{\mathrm{fibre}} = -36$ contributes a cosmological-constant-type term to the effective action, whose role in moduli stabilisation deserves further study.

\item The $24$ non-vanishing shape operator commutators confirm that the gauge sector is genuinely non-abelian, with the commutator structure providing the non-linear part of the field strength.
\end{enumerate}

\paragraph{Relation to standard Kaluza-Klein theory.}
The metric bundle framework differs from standard KK theory~\cite{Bailin1987} in several ways:
(i)~the internal space is non-compact;
(ii)~the section approach replaces the product decomposition;
(iii)~the fibre metric has indefinite signature.
Despite these differences, the mechanism for extracting gauge fields---via the isometry group of the internal space and the Gauss equation---produces the same sign structure as the compact case.

\paragraph{Open questions.}
The fibre isometry group $\SO(4) = \SU(2)_L \times \SU(2)_R$ provides only part of the full Pati-Salam group.
The $\SU(4)$ sector arises from the \emph{structure group} of the normal bundle, not from fibre isometries.
A complete derivation of the $\SU(4)$ gauge kinetic term requires a more detailed analysis of the normal bundle connection, including the role of the complex structure on the Weyl sector~\cite{Paper2}.
The extension from the $\SO(4)$ gauge sector treated here to the full Pati-Salam gauge sector is the subject of ongoing work.
Paper~VI~\cite{Paper6} shows that the indefinite DeWitt metric also constrains observation: the Born rule and uncertainty relations emerge from finite Markov blanket observation of this geometry.

\paragraph{Acknowledgements.}
Computations were performed using NumPy~\cite{numpy} and verified analytically.


% ============================================================
% References
% ============================================================
\bibliographystyle{unsrt}
\bibliography{references}

\end{document}
